\documentclass[10pt,a4paper]{article}
\usepackage{amsmath,amssymb,amsthm,graphicx}
\usepackage{hyperref}
\usepackage{latexsym}
\usepackage{amssymb}
\usepackage{amsthm}
\usepackage{ulem}
\usepackage{enumitem}
\setlist{itemsep=0.4pt,topsep=0pt}
\parindent 0pt
\parindent 0pt
\parskip 12pt
\newcommand{\HCF}{\mathrm{hcf}}
\newcommand{\ZZ}{\mathbb{Z}}
\newcommand{\rk}{\mathrm{rk}}
\DeclareMathOperator{\len}{length}
\DeclareMathOperator{\bl}{b}
\DeclareMathOperator{\cl}{cl}
\DeclareMathOperator{\di}{d}
\DeclareMathOperator{\Cl}{Cl}
\DeclareMathOperator{\acl}{acl}
\DeclareMathOperator{\tp}{tp}
\DeclareMathOperator{\Th}{Th}
\newcommand{\Cor}{\mathcal{C}_\mathrm{or}}
\newcommand{\Cun}{\mathcal{C}_\mathrm{un}}
\newcommand{\mons}{\mathbb{M}}
\newcommand{\strong}{\mathrm{strong}}
\newcommand{\axioms}{\Delta_{\mathrm{FSTS}}}
\newtheoremstyle{note}% hnamei
{8pt}% hSpace abovei
{14pt}% hSpace belowi
{\normalfont}% hBody fonti
{}% hIndent amounti1
{\bfseries}% hTheorem head fonti
{ }% hPunctuation after theorem headi
{.5em}% hSpace after theorem headi2
{}% hTheorem head spec (can be left empty, meaning ‘normal’)i
\theoremstyle{note}
\newtheorem{thm}{Theorem}[section]
%\newtheorem{theorem}{Theorem}
\newtheorem{lemma}[thm]{Lemma} 
\newtheorem{df}[thm]{Definition}
\newtheorem{rmk}[thm]{Remark} 
\newtheorem{prop}[thm]{Proposition}
\newtheorem{fact}[thm]{Fact}
\newtheorem{cor}[thm]{Corollary}
\newtheorem{conj}[thm]{Conjecture}
\newtheorem{ex}[thm]{Example}
\newtheorem{qu}[thm]{Question}
\newtheorem*{notation}{Notation} 
\title{A finite configuration-free partial system has non-negative predimension}
\author{Silvia Barbina, Enrique Casanovas}
\date{arxiv:2411.13723v3; Creative Commons Attribution 4.0 International}
\begin{document}
\maketitle
\noindent\textit{Editorial note.} The counted statement is the author's environment at source0.tex lines 402-402 and its proof environment opens at line 404; the pairing is by adjacency and the retained proof range is the whole environment.\par\medskip
\noindent\textit{Editorial note.} The preamble is the author's own package lines, macro block and theorem-environment declarations (source lines 4-65); the author's optioned amsmath line and the fleqn-dependent layout line are not retained.\par\medskip
\noindent\textit{Editorial note.} The source carries no label-number dump and no counter values are imposed, so the printed numbering is the wrapper's default; identity is the author label, or the pinned source position where the author gave none.\par\medskip
\noindent\textit{Editorial note.} Metadata provenance: the source record of the pinned copy and the licence evidence file are both present; the licence is CC BY 4.0 with status supported, and the card's source\_version matches the value recorded in that evidence file.\par\medskip
\noindent\textit{Editorial scope.} One original result of the pinned author manuscript, counted once: non-negativity of the predimension of a finite configuration-free partial system (unlabelled lemma at source0.tex line 402).\par\medskip
\noindent\textit{Allowed standard prerequisites (editorial).}\begin{itemize}
\item Partial Steiner triple system. A finite structure with a partial operation whose defined products are the blocks.. It is the structure whose predimension is computed.
\item Unconfined configuration. A finite configuration is confined when it forces a product that is not already present; containing none is the hypothesis.. It is what allows an HF-ordering to exist.
\item Predimension. The predimension delta of a finite system is the number of elements minus the number of blocks.. It is the quantity shown to be monotone along the ordering.
\end{itemize}\medskip
\begin{lemma} The predimension of a finite partial STS $A$ that contains no confined configuration is non-negative. \end{lemma}

\begin{proof} 
Suppose $|A|=n$. Since $A$ is unconfined, we can find an HF-ordering %write $A$ as the union of a chain
%\[ \{ a_1 \} \subseteq \{a_1, a_2 \} \subseteq \ldots \subseteq \{a_1, a_2, \ldots, a_n \} \]
\[a_1 < a_2 < \ldots < a_n \, .\]
Let $A_i:= \{a_1, \ldots, a_i \}$ for $1 \leq i \leq n$. Then, if  $|\mathrm{b}(A_i)|$ is the number of blocks in $A_i$,
$$ \delta(A_i)= \begin{cases} |A_{i-1}|- |\mathrm{b}(A_{i-1})| +1 -1 = \delta(A_{i-1}) & \mbox{ if $a_i$ belongs to one block in } A_i \\
|A_{i-1}|- |\mathrm{b}(A_{i-1})| +1 & \mbox{ if $a_i$ belongs to no blocks in } A_i \, .\end{cases} $$
Therefore $\delta(A_i) \geq \delta(A_{i-1})$ for all $i$. 
\end{proof}

\end{document}
